\chapter{Wisdom teeth removal}
\index{Wisdom teeth removal}

\section*{Definition and Overview}
Wisdom teeth removal, or the surgical extraction of the third molars, is one of the most common outpatient surgical procedures performed globally. Wisdom teeth are the final set of permanent teeth to erupt in the human mouth, typically appearing between the ages of 17 and 25 years---a developmental period historically associated with the transition to adulthood and the acquisition of "wisdom." Located at the terminal posterior position of the dental arches (four in total, with one in each quadrant: maxillary right and left, mandibular right and left), these teeth often lack sufficient space within the dental arch for normal eruption and functional alignment.

The clinical significance of wisdom teeth removal stems primarily from the evolutionary mismatch between modern human jaw size and the dental complement. Anthropological evidence suggests that our ancestors possessed larger jaws and coarser diets, which facilitated the accommodation and functional wear of third molars. In contemporary populations, modern dietary patterns, softer foods, and reduced masticatory forces have contributed to a progressive evolutionary reduction in jaw size (specifically the length of the dental arch), without a corresponding reduction in tooth size or number. Consequently, third molars frequently become impacted, meaning they fail to erupt fully into their normal functional position due to physical barriers, such as adjacent teeth or overlying bone and soft tissue.

Impaction can be categorized based on the anatomical position and orientation of the tooth relative to the second molar: mesial (tilted forward toward the front of the mouth), distal (tilted backward toward the rear), vertical (correctly oriented but unable to erupt), and horizontal (lying flat on its side within the jawbone). Impacted third molars can lead to a cascade of local pathological conditions, including pericoronitis (inflammation of the overlying gum tissue), dental caries in both the third molar and the adjacent second molar, periodontal bone loss, root resorption of neighboring teeth, and, in rare instances, the development of odontogenic cysts or tumors. Prophylactic or therapeutic extraction is indicated to treat existing pathology or to prevent future complications, making this procedure a cornerstone of modern oral and maxillofacial surgery.

\section*{Epidemiology}
The epidemiology of third molar impaction and subsequent removal reveals a high prevalence across diverse global populations, with significant implications for public health resources. Research indicates that approximately 65\% to 85\% of young adults exhibit at least one impacted wisdom tooth, with the mandibular third molars being affected slightly more frequently than their maxillary counterparts. The overall prevalence of impaction varies by geographic region and ethnicity, reflecting variations in genetic jaw morphology, dietary habits, and dental development rates. For instance, some studies suggest a lower prevalence of impaction in populations with larger craniofacial skeletons, such as certain Indigenous populations, compared to individuals of European or Asian descent.

The incidence of wisdom teeth removal is highest in individuals aged 18 to 24 years. This peak aligns with the physiological timing of third molar root development and eruption. During this period, the roots are only partially formed (usually one-third to two-thirds of their final length), and the surrounding alveolar bone is relatively elastic, which reduces surgical difficulty and the risk of complications. After age 25, the density and mineralization of the jawbone increase, roots become fully formed and often hook-shaped, and the ligament anchoring the tooth to the bone (periodontal ligament) becomes thinner and less pliable, leading to a higher rate of intraoperative and postoperative complications.

Gender distributions in third molar impaction show a slight female predisposition, with some studies reporting a female-to-male ratio of approximately 1.2:1. This difference is hypothesized to be due to females finishing jaw growth earlier and having slightly smaller dental arches on average than males, reducing the space available for eruption. There is no significant difference in impaction rates between the left and right sides of the mouth, although bilateral impaction is highly common.

\section*{Aetiology and Pathophysiology}
The underlying causes of third molar impaction and the clinical necessity for their extraction are rooted in a combination of evolutionary biology, developmental genetics, and local anatomical factors. The mechanisms leading to pathology include:

\begin{itemize}
    \item \textbf{Evolutionary Jaw Reduction:} The transition from a coarse, abrasive diet of hunter-gatherer societies to the soft, processed diets of modern agricultural and industrial societies has drastically reduced masticatory stimulation. According to the "functional matrix hypothesis," the growth of the mandible and maxilla is stimulated by the forces of chewing. The lack of robust chewing forces in modern humans results in smaller, shorter dental arches that cannot accommodate the full complement of 32 teeth.
    \item \textbf{Delayed Eruption Sequence:} Wisdom teeth are the last teeth to initiate development and erupt. By the time they attempt to enter the oral cavity, the preceding 28 permanent teeth have already occupied the available alveolar space, leaving little to no room for the third molars to assume a functional position.
    \item \textbf{Anatomical Barriers:} Physical obstruction by the adjacent second molar is a primary cause of impaction. When a third molar attempts to erupt at an angle (such as mesioangular), its crown collides with the distal surface of the second molar, halting its progress and trapping it within the bone or soft tissue.
    \item \textbf{Microbial Colonization of the Operculum:} When a third molar partially erupts through the gingiva, it leaves a flap of mucosal tissue (the operculum) overlying the crown. This creates a warm, moist, and anaerobic microenvironment that is virtually impossible to clean with routine oral hygiene. Food debris, desquamated epithelial cells, and plaque accumulate, leading to rapid colonization by anaerobic oral pathogens such as \textit{Prevotella intermedia}, \textit{Fusobacterium nucleatum}, and \textit{Porphyromonas gingivalis}, triggering pericoronitis.
    \item \textbf{Cariogenesis in Hard-to-Reach Areas:} The contact point between a mesioangularly impacted third molar and the adjacent second molar acts as a food trap. Acid-producing bacteria like \textit{Streptococcus mutans} metabolize dietary carbohydrates, producing localized lactic acid that demineralizes the enamel, resulting in dental caries on both the third molar and the distal aspect of the second molar.
    \item \textbf{Pressure-Induced Root Resorption:} In rare cases, the eruptive pathway of the third molar is directed straight into the roots of the second molar. The physical pressure exerted by the erupting crown can trigger osteoclastic activity, leading to pathologic resorption of the second molar's root structure, compromising its structural integrity.
    \item \textbf{Odontogenic Cyst Formation:} The third molar develops within a follicle (a sac of connective tissue and epithelial cells). If the tooth remains impacted, the fluid-filled follicle can undergo cystic degeneration, most commonly developing into a dentigerous cyst. As the cyst expands, it resorbs the surrounding bone and can displace adjacent teeth or, rarely, undergo neoplastic transformation into an ameloblastoma or squamous cell carcinoma.
\end{itemize}

\section*{Clinical Features}
The clinical presentation of wisdom teeth issues ranges from completely asymptomatic, incidentally discovered impactions to acute, severe, life-threatening infections. The signs and symptoms depend on the degree of eruption, position, and presence of secondary infection:

\begin{itemize}
    \item \textbf{Pericoronitis Symptoms:} This is the most common acute presentation. Patients experience localized, throbbing pain in the posterior jaw, which may radiate to the ear, temporomandibular joint (TMJ), or neck. The overlying gum tissue (operculum) appears erythematous, edematous, and is exquisitely tender to palpation. Pus or purulent exudate may be visible draining from under the flap.
    \item \textbf{Trismus:} Difficulty opening the mouth, or trismus, occurs when inflammatory exudate spreads to the adjacent muscles of mastication (primarily the medial pterygoid and masseter muscles). This significantly limits oral opening, hindering speech, mastication, and clinical evaluation.
    \item \textbf{Halitosis and Dysgeusia:} The accumulation of decomposing food debris and anaerobic bacterial activity under the operculum causes a foul, necrotic odor (halitosis) and a persistent bad taste in the mouth (dysgeusia), especially when pressure is applied to the area.
    \item \textbf{Dental Caries Pain:} Chronic food impaction can cause deep decay. Patients may present with thermal sensitivity (pain with hot or cold stimuli), sensitivity to sweet foods, or a dull, lingering ache that indicates irreversible pulpitis or pulpal necrosis.
    \item \textbf{Facial Swelling and Lymphadenopathy:} As infection spreads beyond the alveolar socket, visible swelling of the cheek or angle of the jaw becomes apparent. The submandibular and deep cervical lymph nodes on the affected side become enlarged, firm, and tender to palpation.
    \item \textbf{Systemic Manifestations:} In cases of spreading fascial space infections, patients may exhibit systemic symptoms including fever, malaise, chills, dehydration, tachycardia, and a general feeling of lethargy.
\end{itemize}

\section*{Differential Diagnosis}
When evaluating a patient with posterior jaw pain, it is vital to distinguish third molar pathology from other conditions that mimic its clinical presentation:

\begin{itemize}
    \item \textbf{Temporomandibular Joint Disorder (TMD):} Characterized by pain in the jaw joint and masticatory muscles, TMJ clicking or popping, and limited range of motion. Unlike pericoronitis, TMD pain is exacerbated by jaw function, is not associated with localized gingival swelling or purulence, and shows no radiographic abnormalities of the third molars.
    \item \textbf{Dental Caries or Pulpal Pathology of the First or Second Molars:} Pain originating from deep decay or pulpal infection in the adjacent molars can mimic third molar pain. Detailed clinical examinations, including percussion testing, thermal testing, and periapical radiographs, help isolate the affected tooth.
    \item \textbf{Periodontitis:} Generalized or localized periodontal disease can cause bone loss, gum swelling, and pain in the posterior regions. However, periodontitis is typically widespread, characterized by deep periodontal pocketing across multiple teeth, and lacks the specific anatomical involvement of an erupting operculum.
    \item \textbf{Myofascial Pain Syndrome:} Chronic muscular pain involving trigger points in the masseter, temporalis, or pterygoid muscles. The pain is dull and aching, often associated with bruxism (teeth grinding) or stress, and lacks signs of acute infection or dental pathology.
    \item \textbf{Sinusitis:} Maxillary sinusitis can cause referred pain to the upper teeth, particularly the maxillary molars, due to the close anatomical proximity of the tooth roots to the maxillary sinus floor. Patients with sinusitis typically report nasal congestion, rhinorrhea, facial pressure that worsens when bending forward, and bilateral tenderness of multiple upper teeth, rather than localized pain in the lower jaw or isolated to one third molar.
    \item \textbf{Salivary Gland Pathology:} Conditions affecting the submandibular or parotid glands, such as sialolithiasis (salivary stones) or sialadenitis (gland infection), can cause pain and swelling near the angle of the jaw. Pain is typically exacerbated by eating (due to salivary stimulation) and can be differentiated by examining salivary flow and palpating the gland ducts.
    \item \textbf{Trigeminal Neuralgia:} A neurological condition characterized by sudden, severe, lancinating, shock-like pain triggered by light touch to the face. The absence of dental pathology on radiographs and the classic trigger-zone presentation differentiate it from dental pain.
\end{itemize}

\section*{When to Seek Medical Care}
Patients experiencing signs or symptoms of wisdom teeth complications should consult a dentist or oral surgeon. Prompt evaluation is recommended if there is persistent jaw pain, bleeding gums, difficulty opening the mouth, or swelling.

Immediate emergency medical attention should be sought if the patient develops signs of a spreading, life-threatening infection. Go to the nearest emergency department or call emergency services---such as \textbf{local emergency services} in many countries, \textbf{911} in the United States, or \textbf{999} in the United Kingdom---if the following red-flag symptoms occur:
\begin{itemize}
    \item Swelling of the neck, submental space, or floor of the mouth that causes difficulty swallowing (dysphagia) or difficulty breathing (dyspnea).
    \item A rapidly expanding swelling of the face or eye area, indicating possible orbital or cavernous sinus involvement.
    \item A high fever (above $38.5\,^{\circ}C$ or $101.3\,^{\circ}F$) accompanied by rigors, confusion, or severe lethargy.
    \item Inability to swallow fluids or control secretions, leading to drooling (sialorrhea), which suggests impending airway obstruction.
\end{itemize}

\section*{Diagnosis and Investigations}
A definitive diagnosis and treatment plan for wisdom teeth removal require a comprehensive clinical examination combined with specialized dental imaging:

\begin{itemize}
    \item \textbf{Clinical Extraoral Examination:} The clinician assesses facial symmetry, palpates the temporomandibular joints, evaluates the range of mandibular motion, and checks for enlarged, tender submandibular and cervical lymph nodes.
    \item \textbf{Clinical Intraoral Examination:} The oral cavity is inspected for the presence of wisdom teeth, mucosal swelling, erythema, operculum status, purulence, and dental hygiene. Probing is performed around the second and third molars to detect periodontal pockets. Percussion and vitality tests are completed on adjacent teeth to rule out pulpal disease.
    \item \textbf{Panoramic Radiograph (Orthopantomogram or OPG):} This is the gold standard imaging modality for initial assessment. It provides a broad, two-dimensional view of both the maxilla and mandible, showing the position, angulation, root morphology, and degree of impaction of all four third molars. Importantly, it allows the surgeon to assess the relationship between the mandibular third molar roots and the inferior alveolar nerve (IAN) canal, as well as the proximity of maxillary third molar roots to the maxillary sinus.
    \item \textbf{Cone Beam Computed Tomography (CBCT):} When a panoramic radiograph shows high-risk features (such as darkening of the tooth root, diversion of the IAN canal, or narrowing of the root where it crosses the canal), a CBCT is indicated. CBCT provides high-resolution, three-dimensional reconstruction of the jawbone, enabling the surgeon to precisely map the anatomical course of the IAN relative to the roots (e.g., lingual, buccal, or inter-radicular), significantly reducing the risk of nerve injury during surgery.
    \item \textbf{Periapical Radiographs:} Occasionally used for detailed, high-resolution views of the crown or root apex of a specific tooth, although they are often difficult to position comfortbly in patients with severe trismus or gag reflexes.
\end{itemize}

\section*{Management}
The management of third molar issues ranges from conservative observation to surgical extraction. The choice of treatment depends on the presence of symptoms, pathology, patient age, and surgical risk:

\begin{itemize}
    \item \textbf{Active Surveillance:} For asymptomatic, disease-free, fully impacted teeth that are deeply embedded in bone with no radiographic signs of pathology, active surveillance is often recommended. This involves periodic clinical and radiographic monitoring (every 1 to 2 years) to detect early signs of decay, bone loss, or cyst formation without subjecting the patient to the risks of surgery.
    \item \textbf{Pharmacotherapy for Acute Infection:} In cases of active pericoronitis, local and systemic medical management is initiated before surgery. Local management includes irrigating under the operculum with sterile saline or a chlorhexidine gluconate solution to remove debris. Systemic antibiotics (such as amoxicillin, metronidazole, or phenoxymethylpenicillin) are prescribed if the patient exhibits systemic symptoms, spreading swelling, or trismus. Nonsteroidal anti-inflammatory drugs (NSAIDs) like ibuprofen, or paracetamol (acetaminophen), are recommended for pain control.
    \item \textbf{Surgical Extraction:} This is the definitive treatment for symptomatic or pathological third molars. The procedure is typically performed under local anesthesia, local anesthesia with intravenous (IV) conscious sedation (ideal for anxious patients), or general anesthesia (for complex, deeply impacted cases or multi-quadrant surgeries).
    \item \textbf{Surgical Technique:} The standard surgical steps involve:
    \begin{enumerate}
        \item Creating a mucoperiosteal flap by making an incision in the gingiva to expose the underlying alveolar bone and tooth.
        \item Osteotomy (bone removal) using a surgical handpiece and bur to create space and expose the crown of the impacted tooth.
        \item Odontosection (tooth sectioning), where the tooth is cut into pieces (e.g., separating the crown from the roots or splitting the roots) to allow extraction through a smaller bony window, preserving surrounding bone.
        \item Elevation and luxation of the tooth segments using specialized instruments (elevators and forceps).
        \item Debridement of the socket, including removal of the remaining dental follicle and smoothing of sharp bony edges.
        \item Closure of the wound using sutures (often resorbable, such as chromic gut or vicryl).
    \end{enumerate}
    \item \textbf{Coronectomy (Intentional Root Retention):} When a mandibular third molar is deeply impacted and radiographic imaging (CBCT) shows that its roots directly envelope or compress the inferior alveolar nerve, a coronectomy may be performed. In this procedure, the crown of the tooth is surgically removed, but the roots are intentionally left in place. The pulp chamber is treated, and the mucosal flap is closed over the roots. This technique has been shown to reduce the risk of permanent nerve damage to less than 1\%, with a low rate of subsequent root migration or infection.
\end{itemize}

\section*{Complications}
While wisdom teeth removal is generally safe, it is associated with several well-documented complications, ranging from minor, self-limiting issues to permanent neurological deficits:

\begin{itemize}
    \item \textbf{Alveolar Osteitis (Dry Socket):} Occurs in approximately 5\% to 15\% of extractions, most commonly in the mandible. It happens when the blood clot in the socket fails to form, dislodges, or premature dissolves (often due to fibrinolysis). This exposes the underlying alveolar bone, causing severe, throbbing pain that radiates to the ear, typically starting 3 to 4 days postoperatively. Risk factors include smoking, oral contraceptive use, poor oral hygiene, and traumatic extractions.
    \item \textbf{Sensory Nerve Injury:} The inferior alveolar nerve (IAN) and the lingual nerve run in close proximity to the mandibular third molars. Damage to these nerves during extraction (due to traction, compression, or transection) can result in paresthesis (numbness), dysesthesia (altered sensation), or hypoesthesia (reduced sensation) of the lower lip, chin, teeth, gums (for IAN), or the anterior two-thirds of the tongue and taste sensation (for the lingual nerve). While most nerve injuries resolve within 3 to 6 months, approximately 0.5\% to 1\% result in permanent sensory deficits.
    \item \textbf{Maxillary Sinus Communication (Oroantral Fistula):} The roots of the maxillary third molars lie near the floor of the maxillary sinus. Extraction can lead to a perforation of the thin bone separating the socket from the sinus, creating an oroantral communication. Small openings ($<2$\,mm) usually heal spontaneously, but larger openings require surgical closure to prevent chronic sinusitis and fluid passage from the mouth into the nose.
    \item \textbf{Postoperative Infection and Abscess:} Occurs in 1\% to 4\% of cases. It presents as increasing pain, swelling, and purulence 1 to 2 weeks after surgery, requiring wound debridement, irrigation, and antibiotics.
    \item \textbf{Mandibular Fracture:} An extremely rare but serious complication ($<0.1$\% of cases) that occurs either intraoperatively due to excessive force applied with elevators, or postoperatively (typically within 2 to 4 weeks) due to weakened jaw structure combined with accidental trauma or chewing hard foods.
    \item \textbf{Damage to Adjacent Teeth:} The surgical instruments or the extraction path can cause accidental damage, such as chipping, luxation, or crown fracture, to the adjacent second molars.
\end{itemize}

\section*{Prognosis and Outcomes}
The overall prognosis for patients undergoing wisdom teeth removal is excellent. For the vast majority of patients, the procedure resolves the underlying pathology, prevents future infections, and eliminates pain.

The recovery phase typically follows a predictable course:
\begin{itemize}
    \item \textbf{Days 1 to 2:} The acute inflammatory phase. Patients experience peak swelling and discomfort. Mild bleeding or oozing from the sockets is normal and expected.
    \item \textbf{Days 3 to 5:} Swelling begins to subside, and oral opening starts to improve. Pain transition from moderate to mild, requiring less analgesic medication.
    \item \textbf{Days 7 to 10:} Mucosal healing is largely complete, and sutures (if resorbable) begin to fall out or dissolve. Most patients can return to their normal diet and physical activities.
    \item \textbf{Weeks 4 to 8:} Complete bone remodeling within the extraction socket occurs, filling the void with new trabecular bone.
\end{itemize}

The factors that adversely affect prognosis and increase recovery time include advanced age (over 25 years), deep bony impactions, history of smoking, systemic medical comorbidities (such as poorly controlled diabetes or immunosuppression), and poor adherence to postoperative care instructions.

\section*{Prevention and Self-Care}
While the developmental impaction of wisdom teeth cannot be prevented, complications can be minimized through early clinical screening, routine dental visits, and meticulous postoperative self-care:

\begin{itemize}
    \item \textbf{Early Screening and Monitoring:} Regular dental checkups during adolescence (ages 14 to 18) with panoramic imaging allow for early detection of potential impactions. Extraction can be planned proactively before root development is complete, minimizing surgical risks.
    \item \textbf{Immediate Postoperative Care (First 24 Hours):}
    \begin{itemize}
        \item Apply pressure by biting on a sterile gauze pad for 30 to 45 minutes to promote stable clot formation.
        \item Apply ice packs to the side of the face in 20-minute intervals to reduce facial swelling.
        \item Avoid spitting, sucking through a straw, smoking, or drinking carbonated and hot beverages, as the negative pressure and heat can dislodge the blood clot.
        \item Consume a soft, cool diet (e.g., yogurt, smoothies, pudding, applesauce).
    \end{itemize}
    \item \textbf{Subsequent Postoperative Care (Day 2 Onward):}
    \begin{itemize}
        \item Gently rinse the mouth with warm salt water (1/2 teaspoon of salt in a glass of warm water) 4 to 5 times a day, especially after meals, to keep the sockets clean.
        \item Resume gentle brushing of the teeth, avoiding the extraction sites, to maintain overall oral hygiene.
        \item Keep hydrated and gradually reintroduce solid foods as pain and jaw opening permit.
        \item Avoid strenuous physical exercise for the first 3 to 5 days, as elevated blood pressure can cause secondary bleeding from the sockets.
    \end{itemize}
\end{itemize}

\section*{Sources and Further Reading}
\begin{itemize}
    \item \textbf{American Association of Oral and Maxillofacial Surgeons (AAOMS):} Clinical practice guidelines and patient resources on third molar management. \url{https://www.aaoms.org}
    \item \textbf{National Institute for Health and Care Excellence (NICE):} UK clinical guidelines on the extraction of wisdom teeth. \url{https://www.nice.org.uk}
    \item \textbf{Mayo Clinic:} Comprehensive patient information on wisdom tooth extraction procedure and recovery. \url{https://www.mayoclinic.org}
    \item \textbf{Australian Dental Association (ADA):} Patient education materials on oral surgery and third molar management. \url{https://www.ada.org.au}
    \item \textbf{British Association of Oral and Maxillofacial Surgeons (BAOMS):} Professional guidelines and patient information leaflets on wisdom teeth removal. \url{https://www.baoms.org.uk}
    \item \textbf{Cochrane Database of Systematic Reviews:} Evidence-based reviews on prophylactic extraction versus surveillance of third molars. \url{https://www.cochranelibrary.com}
\end{itemize}